\section{The datum, the signed sum and the two lattice bounds}\label{sz:setup}
For a positive multiple $q_0$ of $18480$ and a prime $d\equiv q_0/2-1\pmod{q_0}$,
write $P=(q_0,d)$. The following conditions specify an admissible datum
$D$ for $P$. Theorem~\ref{thm:p1-data} constructs such data at every
sufficiently large integer length. The construction hypotheses on forms
and lifts will be imposed separately for each datum.

\paragraph{Scales and places.}
The positive integers $x,x_1,b$ have logarithms satisfying
\[
 10^{12}C_P\le\log L,\quad L^{100}\le x,\quad x^{10^6}\le x_1,
 \quad x_1^{10^6}\le b,\quad (\log x_1)^6\le L,\qquad L=\log b .
\]
The parameter constant $C_P$ is the fixed constant in
\srcline{New/Reduction/Data/AdmissibleData}{50}{the definition of admissibility}.
In particular $C_P\ge64$ and it bounds
$|c_K|$, where $c_K=\gamma-\Re(L'/L)(1,\chi_d)$ and $\chi_d(n)=(n/d)$.
The length $m$ is a positive power of two in $[\sqrt x,2\sqrt x]$ and the
unitary rank $r$ a positive power of three in $[x^5,3x^5]$.
Take $s=22\log_2 m\le32\log x$,
$\ell_0=3^{\lfloor\log_3s\rfloor+4}$,
$A=3\ell_0 2^s$, $q_1=q_0A$ and $q_2=(q_0d)^{s+6}$, with $q_1x_1\mid b$.
Put $a=xb$, $K=\mathbb Q(\sqrt{-d})$ and $E=K(\zeta_{q_1})$.
Represent its infinite places by the embeddings extending a fixed embedding of
$K$. They form the group $\Phi\simeq(\mathbb Z/q_1)^\times$ and $f=|\Phi|$.
The fibre $\Phi_a$ over $1\bmod q_0$ has size $A$. In the fibre over
$-1$ choose $\Phi_b$ with $|\Phi_b|L\ge A$. Put $k=r+1-m$.
For every prime $p\le k(5r)$ and $0<i<j<\ell_0$, there are no
$\sigma,\tau,\rho\in\Phi_b$ with
$\tau/\sigma\equiv\pm p^i$, $\rho/\sigma\equiv\pm p^j$
modulo the prime-to-$p$ part of $q_1$, with the signs chosen independently.
The place weight $\kappa(\tau)$ equals $b$ on $\Phi_b$ and $a/x_1$ elsewhere.

The distinct level primes $p_0,p_1$ lie in $[b^{10},b^{11}]$ and the
auxiliary prime $N$ has the logarithmic window below. All three are
$1\bmod q_2^2$. Write $M=p_0p_1$, $\log\bar\kappa=f^{-1}\sum_\tau\log\kappa(\tau)$,
and $\mathrm{rd}_E=|\mathrm{disc}(E)|^{1/(2f)}$. There is an integer of $K$
whose product with its conjugate is $M$ and an integer of $E$ with corresponding
product $N$; these choices are part of $D$. The required window is
\[
 \log N_-\le\log N\le\log N_-+\sqrt L,\quad
 \log N_-=m\log\bar\kappa+(m-1)\log(M\mathrm{rd}_E/(2\pi))-15L/2 .
\]

\paragraph{The two scalar quantities.}
For $h\in\mathbb N$ put $j(h)=\log(h!)+h\log(M/\Omega)$,
where $\Omega=2\pi e^{c_K}/\mathrm{rd}_E$. Set $u(t)=t\bmod q_1$ if $t$ is a
unit modulo $q_1$, and $u(t)=1$ otherwise. With $v_+=\max(v,0)$ set
\[
 S(y,\tau)=\sum_{1\le\nu\le a}\Lambda(\nu)(1-\chi_d(\nu))
       (y/\nu-\kappa(u(-\nu)\tau))_+,
 \qquad S_0=a(\log a-\log\bar\kappa+c_1).
\]
Here $c_1=c_{d,1}$ is the modulus-one centering constant in \eqref{eq:R}. Explicitly
$c_1=\Re[(L'/L)(1,\chi_d)-((s-1)\zeta(s))'/((s-1)\zeta(s))|_{s=1}-1]$,
where $(s-1)\zeta(s)$ has its regular value at $1$. Thus $c_1=-c_K-1$.
With the normalization of \eqref{eq:total}, the signed excess
and the common main height are
\begin{equation}\label{sz:height}
 \Tot_D=\sum_{\tau\in\Phi_a}(S(a,\tau)-S_0),\qquad
 H=mk\sum_{\tau\in\Phi}\bigl(j(\kappa(\tau))-S(\kappa(\tau),\tau)\bigr).
\end{equation}
Write $U=faL^3$ and $V=faL^4$. The arguments of $S$ and the sets of places
in these two sums differ. Both determinant estimates use exactly this $H$.

\paragraph{The spaces and maps.}
Let $X$ be the space of algebraic forms on the definite rank-$k$ space
$X_0$ at its prescribed level and Weyl weight. At each of the $f$ chosen places
$w\mid N$ there is a central scaling translation $T_w$ of order dividing
$N-1$. The ramified source is $S=\bigcap_w\ker(\sum_{i=0}^{N-2}T_w^i)$.
The argument of $S(y,\tau)$ distinguishes the prime sum from this space.
Write $Y$ for the principal-level automorphic forms of the prescribed type
on the rank-$r$ space $Y_0$, with the Petersson inner product. The finite
residue group $G$ at the chosen places of $p_0,p_1$ acts through $\rho$.
The central good-type projector $e_0$ has image $Y_T=e_0Y$.
The domain dimension of the normalized maps is $d_S=\dim_\mathbb C S$,
their target dimension is $d_T=\dim_\mathbb C Y_T$, and $d_X=\dim_\mathbb C X$.
For the Poincar\'e map $\Pi:X\longrightarrow Y$ defined by the chosen Bessel
functional and point, the normalization used throughout is
\[
 P_g=c_B c_{\rm ss}^{-1}|G|\,e_0\rho(g)\Pi|_S,\qquad g\in G .
\]
The seesaw gives the common unit $c_{\rm ss}$, and the local divisibility
bounds determine the positive rational $c_B$ in Section~\ref{A:section}.
Centrality implies $P_g=\rho(g)P_1$. At each infinite place $Y_0$ has signature $(r-m/2,m/2)$: the initial
string of length $m-1$ contributes $m/2-1$ hyperbolic pairs and one
negative central line, in addition to $k$ positive directions. The chain $\Theta_\xi$ takes its good forms through $m$ lifts
to algebraic forms on the definite terminal space $Y_m$.

An $m$-fold twist $\xi$ is a weight-zero tame automorphic character of
$U(1)$ with the admissible finite support disjoint from its conjugate
and of norm prime to $MN$. Every power $\xi^j$, $1\le j\le m$, has
nontrivial tame character at every place of that support. Let $N_\xi$ be
the product of the residue norms on the support. It is allowed at $p$ if
$p\nmid N_\xi$. The empty-support trivial twist is denoted by $1$.
For an identification $\iota:\mathbb C\simeq\overline{\mathbb Q}_p$ write
$|z|_\iota=|\iota z|_p$ with $|p|_p=p^{-1}$ and $O_\iota=\{z:|z|_\iota\le1\}$.
Integrality of an algebraic form means that each adelic value, transported
back by the inverse finite local action, lies in the integral Weyl module.
Write $\mathcal I\subset S$ for source forms integral for every $p,\iota$
and $\mathcal L_{\xi,\iota}$ for the corresponding module at the end of the chain.

\paragraph{Index and squared covolume.}
For a finite indexed family $y=(y_i)$ in $Y_T$, the assertion
$\operatorname{IndexLe}(y,\Delta)$ has the following order of quantifiers.
There exist numbers $e_p\ge0$ with $\sum_{p\in F}e_p\log p\le\Delta$ for
every finite set of primes $F$. For every prime $p$ and every $\iota$ above it
there exists one square matrix $u_\iota$ over $O_\iota$ on the index set of
$y$, satisfying
\begin{equation}\label{sz:index}
 |\det u_\iota|_\iota\ge p^{-e_p},\qquad
 \Theta_\xi\!\left(\sum_j (u_\iota)_{ij}y_j\right)\in\mathcal L_{\xi,\iota}
 \quad\text{for all }i\text{ and all }\xi\text{ allowed at }p.
\end{equation}
At a fixed identification the same rows must work for every allowed twist.

Let $v_\sigma(z)$ be the vector obtained by applying
$\sigma\in\operatorname{Aut}(\mathbb C)$ to each of the prescribed
coordinates of $\Theta_1z$ at the terminal class representatives.
The class pairing has a rational positive definite matrix $G_{\rm cl}$
in these coordinates. Set $N_\sigma(z)=v_\sigma(z)^*G_{\rm cl}v_\sigma(z)$.
Then
\begin{equation}\label{sz:covol}
 \operatorname{CovolumeLe}(y,c)\quad\Longleftrightarrow\quad
 \left|\det\bigl(v_\sigma(y_i)^*G_{\rm cl}v_\sigma(y_j)\bigr)_{ij}\right|
       \le e^c\qquad(\sigma\in\operatorname{Aut}(\mathbb C)).
\end{equation}
The Gram matrix is formed \emph{after} acting on coordinates, so it is positive
semidefinite for every field automorphism. This order matters even with rational
$G_{\rm cl}$, since an arbitrary $\sigma$ can fail to commute with complex conjugation.

\paragraph{Statement and dependencies.}
For a fixed admissible $D$, the hypotheses of
Sections~\ref{A:section}, \ref{B:frontier} and \ref{c3:frontier}
imply the pointwise determinant estimate \eqref{eq:det}.
Theorems A and B produce an independent family with small squared
covolume and controlled index. Theorem C applies to that same family.
Section~\ref{sz:comparison} deduces the estimate and gives the
quantifiers required for Hypothesis~\ref{hyp:scale}.
Appendix~\ref{sz:register} records the construction steps that
remain hypotheses.
